\section*{Chapter \ref{ch:b1:alcohol-activation} --- Alcohols: Activating the OH Group}

\begin{solution}{exo:b1:alcohol-activation:1}
\ce{2CH3CH2OH + 2Na -> 2CH3CH2O- + 2Na+ + H2};
\ce{CH3CH2OH + NaH -> CH3CH2O- + Na+ + H2};
\ce{CH3CH2OH + OH- <=> CH3CH2O- + H2O}, $K = 10^{14.0 - 15.9} = 10^{-1.9}$:
partial.
\end{solution}

\begin{solution}{exo:b1:alcohol-activation:2}
1-Methoxypropane; methoxyethane; ethoxybenzene (the phenoxide, a weaker
base than an \omterm{def:b1:alcohol-activation:alkoxide}{alkoxide}, is still a good \omterm{def:b1:electronic-effects:nucleophile}{nucleophile}).
\end{solution}

\begin{solution}{exo:b1:alcohol-activation:3}
Tosylation in pyridine gives \ce{CH3CH2CH2OTs}; then
\[ \ce{CH3CH2CH2OTs + I- -> CH3CH2CH2I + TsO-}. \] Directly, iodide would have to
expel \ce{OH-}: no reaction. The \omterm{def:b1:alcohol-activation:sulfonate}{tosylate} has an excellent \omterm{def:b1:electronic-effects:nucleophile}{leaving group},
and the conditions are neither acidic nor strongly basic.
\end{solution}

\begin{solution}{exo:b1:alcohol-activation:4}
Butan-2-ol: but-1-ene, \cip{Z}- and \cip{E}-but-2-ene; major
\cip{E}-but-2-ene. 2-Methylpentan-2-ol: 2-methylpent-2-ene (major,
trisubstituted) and 2-methylpent-1-ene.
\end{solution}

\begin{solution}{exo:b1:alcohol-activation:5}
\ce{(CH3)2CHCH2-O-CH2CH3}: sodium 2-methylpropan-1-olate with bromoethane,
or sodium ethoxide with 1-bromo-2-methylpropane. Both halides are primary,
but 1-bromo-2-methylpropane is branched next to the reacting carbon,
which slows SN2 and lets the \omterm{def:b1:predominant-reaction:strong-acid}{strong base} eliminate (2-methylpropene). The
first choice is better.
\end{solution}

\begin{solution}{exo:b1:alcohol-activation:6}
The \omterm{def:b1:alcohol-activation:sulfonate}{tosylate} keeps the \cip{R} \omterm{def:b1:stereochemistry-in-depth:isomers}{configuration}; SN2 with ethoxide inverts
it: \cip{S}-2-ethoxyoctane (OEt ranks first as OTs did). With sulfuric acid
and heat: elimination to octenes (mainly \cip{E}-oct-2-ene), via a
secondary \omterm{def:b1:electronic-effects:intermediates}{carbocation}, with loss of the stereochemistry.
\end{solution}

\begin{solution}{exo:b1:alcohol-activation:7}
Protonation of OH; loss of water to the tertiary \omterm{def:b1:electronic-effects:intermediates}{carbocation} (slow);
capture by \ce{Cl-}: 2-chloro-2-methylpropane. Fast because the tertiary
\omterm{def:b1:electronic-effects:intermediates}{carbocation} is stable and the chloride, insoluble in the acid, separates.
\end{solution}

\begin{solution}{exo:b1:alcohol-activation:8}
\ce{CH3CH2OH + H+ -> CH3CH2OH2+}; a second ethanol attacks the carbon
from behind, water leaving (SN2): \ce{(CH3CH2)2OH+}, which loses a proton:
ethoxyethane. Competing: a base (ethanol, \ce{HSO4-}) takes a $\beta$
proton of \ce{CH3CH2OH2+} as water leaves (E2): ethene.
\end{solution}

\begin{solution}{exo:b1:alcohol-activation:9}
2-Methylpropan-2-ol (tertiary), giving 2-chloro-2-methylpropane, insoluble.
\end{solution}

\begin{solution}{exo:b1:alcohol-activation:10}
Protonation and loss of water give a secondary \omterm{def:b1:electronic-effects:intermediates}{carbocation} next to a
quaternary carbon. A methyl group of that carbon moves with its \omterm{def:b1:lewis-resonance-vsepr:lewis}{bonding
pair} to the cationic carbon: the positive charge is now on a \omterm{def:b1:electronic-effects:carbon-class}{tertiary
carbon}, more stable. Loss of a proton from a neighbouring carbon gives the
tetrasubstituted 2,3-dimethylbut-2-ene, the most stable alkene.
\end{solution}

\begin{solution}{exo:b1:alcohol-activation:11}
From \cip{R}-butan-2-ol: \omterm{def:b1:alcohol-activation:sulfonate}{tosylate} (\cip{R}), then sodium methoxide (SN2,
inversion): \cip{S}-2-methoxybutane. From \cip{S}-butan-2-ol: make the
\omterm{def:b1:alcohol-activation:alkoxide}{alkoxide} (NaH) and add iodomethane: the stereogenic C--O bond is not
touched, \cip{S}-2-methoxybutane again.
\end{solution}

\begin{solution}{exo:b1:alcohol-activation:12}
$Q = [\text{ether}][\ce{H2O}]/[\ce{EtOH}]^2$. Removing water as it forms
keeps $[\ce{H2O}]$ and $Q$ small: $Q < K$ and the reaction goes on forming
ether until the alcohol is consumed.
\end{solution}

\begin{solution}{pb:b1:alcohol-activation:1}
\textbf{1.} Sodium tert-butoxide with a methyl halide; sodium methoxide with
a tert-butyl halide.
\textbf{2.} The second: methoxide, a \omterm{def:b1:predominant-reaction:strong-acid}{strong base}, eliminates from the
tertiary halide, \ce{(CH3)3CBr + CH3O- -> (CH3)2C=CH2 + CH3OH + Br-}.
\textbf{3.} \ce{(CH3)3CO- + CH3I -> (CH3)3COCH3 + I-}: the \omterm{def:b1:alcohol-activation:alkoxide}{alkoxide} oxygen
attacks the methyl carbon from behind as iodide leaves (SN2).
\textbf{4.} With sodium (or sodium hydride): \ce{2(CH3)3COH + 2Na ->
2(CH3)3CO- + 2Na+ + H2}. Hydroxide is too weak a base: $K = 10^{14.0 -
19.2} = 10^{-5.2}$.
\textbf{5.} Water would protonate the \omterm{def:b1:alcohol-activation:alkoxide}{alkoxide} and attack the halide.
\textbf{6.} No: no \omterm{def:b1:stereochemistry-in-depth:stereocentre}{stereogenic centre}.
\textbf{7.} \cip{R}-butan-2-ol + TsCl (pyridine) $\to$ \cip{R}-butan-2-yl
\omterm{def:b1:alcohol-activation:sulfonate}{tosylate}: retention.
\textbf{8.} \cip{S}-2-methoxybutane, by SN2 inversion.
\textbf{9.} \cip{S}-butan-2-ol.
\textbf{10.} \cip{R}-2-methoxybutane: the C--O bond of the stereogenic
carbon is kept; only O--C(methyl) forms.
\textbf{11.} E2 (methoxide is a \omterm{def:b1:predominant-reaction:strong-acid}{strong base} and the carbon is secondary),
giving butenes; limited by the small base, an excellent \omterm{def:b1:electronic-effects:nucleophile}{leaving group} and
mild temperature.
\textbf{12.} The protonated alcohol ionises to a planar secondary
\omterm{def:b1:electronic-effects:intermediates}{carbocation}, captured on both faces by methanol: racemic ether, together
with butenes.
\textbf{13.} Protonation of OH, loss of water (slow) to the tertiary
\omterm{def:b1:electronic-effects:intermediates}{carbocation}, loss of a $\beta$ proton.
\textbf{14.} 2-Methylbut-2-ene (major, trisubstituted) and
2-methylbut-1-ene.
\textbf{15.} A tertiary \omterm{def:b1:electronic-effects:intermediates}{carbocation} is hindered and loses a proton far
more easily than another alcohol molecule can attack it.
\textbf{16.} To remove a product: with $Q$ kept small the reaction goes on,
and the alkene is not left in contact with the acid.
\textbf{17.} A high first barrier (loss of water, rate-determining) to the
\omterm{def:b1:electronic-effects:intermediates}{carbocation}, then a small barrier to the alkene.
\textbf{18.} No: a primary \omterm{def:b1:electronic-effects:intermediates}{carbocation} is too unstable; primary alcohols
eliminate by E2 on the protonated alcohol, at higher temperature.
\textbf{19.} $10.0/74.0 = \qty{0.135}{mol}$.
\textbf{20.} $1.10 \times 0.135 \times 141.9 = \qty{21.1}{g}$.
\textbf{21.} $0.135 \times 88.0 = \qty{11.9}{g}$.
\textbf{22.} $0.75 \times 11.9 =$ \textbf{\qty{8.9}{g} of MTBE}.
\end{solution}
